% Generated by scripts/assemble.py; corrections belong in source/corrections.json.
\examhead{142}{湖南大学 2025 年数学分析试卷}{135}
% p136-b20; source PDF page 136
\blockanchor{p136-b20}
\question{1}已知 $x_{1} = 1, x_{n + 1} = \sqrt{5x_{n}}$，证明：$\{x_{n}\}$ 收敛并求其极限.\par
% p136-b21; source PDF page 136
\blockanchor{p136-b21}
\question{2}已知 $\{a_{n}\}$ 是无界数列, 但不是无穷大量. 证明: $\{a_{n}\}$ 存在两个子列, 一个是收敛数列, 另一个是无穷大量.\par
% p136-b22; source PDF page 136
\blockanchor{p136-b22}
\question{3}设 $f(x)$ 是 $\mathbb{R}$ 上的连续函数, 若 $\lim_{|x|\to +\infty}f(x) = +\infty$, 则存在 $x_0 \in \mathbb{R}$, 使得 $f(x)$ 可以取到 $\inf_{x \in \mathbb{R}} f(x)$.\par
% p136-b23; source PDF page 136
\blockanchor{p136-b23}
\question{4}解答如下问题:\par
% p136-b24; source PDF page 136
\blockanchor{p136-b24}
(1) 若 $0 < \eta < 1$, 证明: $f(x) = \sin \frac{1}{x^2}$ 在 $(\eta, 1)$ 上一致连续.\par
% p136-b25; source PDF page 136
\blockanchor{p136-b25}
(2) 证明 $f(x) = \sin \frac{1}{x^2}$ 在 (0,1) 上不一致连续.\par
% p137-b00; source PDF page 137
\blockanchor{p137-b00}
\question{5}讨论 $\sum_{n=1}^{\infty} \frac{\sin nx}{n^p} (p \in \mathbb{R})$ 的敛散性.\par
% p137-b01; source PDF page 137
\blockanchor{p137-b01}
\question{6}设 $f(x)$ 在 $(a, b)$ 上可导, 任取 $x_1, x_2 \in (a, b)$, 满足 $f'(x_1)f'(x_2) < 0$. 证明: 至少存在一点 $\xi$ 介于 $x_1, x_2$ 之间, 使得 $f'(\xi) = 0$.\par
% p137-b02; source PDF page 137
\blockanchor{p137-b02}
\question{7}解答如下问题:\par
% p137-b03; source PDF page 137
\blockanchor{p137-b03}
(1) 证明: $\int_{a}^{b}\left|x - \frac{a + b}{2}\right|^{n}\mathrm{d}x = \frac{b^{n + 1} - a^{n + 1}}{2^n(n + 1)}.$\par
% p137-b04; source PDF page 137
\par\Needspace{6\baselineskip}
\blockanchor{p137-b04}
(2) 已知  $\int_{0}^{1} x^{n} f(x) \, dx = 1, \int_{0}^{1} x^{k} f(x) \, dx = 0 (k = 0, 1, 2, \cdots, n - 1)$ . 证明:\par
\nopagebreak[4]
% p137-b05; source PDF page 137
\blockanchor{p137-b05}
\displayformula{\max _ {x \in [ 0, 1 ]} | f (x) | \geq 2 ^ {n} (n + 1).}
% p137-b06; source PDF page 137
\blockanchor{p137-b06}
\question{8}证明: $x \to a$ 时 $f(x, y)$ 一致收敛于 $\varphi(y)$ 的充要条件是对任意趋近于 $a$ 的数列 $\{x_n\}$, $f(x_n, y)$ 一致收敛于 $\varphi(y)$.\par
% p137-b07; source PDF page 137
\par\Needspace{6\baselineskip}
\blockanchor{p137-b07}
\question{9}计算\par
\nopagebreak[4]
% p137-b08; source PDF page 137
\blockanchor{p137-b08}
\displayformula{I = \iint_ {\Sigma} 4 x z \mathrm{d} y \mathrm{d} z - 2 y z \mathrm{d} z \mathrm{d} x + (1 - z ^ {2}) \mathrm{d} x \mathrm{d} y.}
% p137-b09; source PDF page 137
\blockanchor{p137-b09}
其中  $\Sigma$  是由  $z = e^{y} (0 \leq y \leq a)$  绕 z 轴旋转一周生成的曲面, 取下侧.\par
